% Generated draft: requires mathematical and source review.
\documentclass[11pt]{article}
\usepackage[margin=25mm]{geometry}
\usepackage{fontspec}
\setmainfont{Latin Modern Roman}
\usepackage{amsmath,amssymb,mathtools,unicode-math}
\setmathfont{Latin Modern Math}
\usepackage{enumitem,needspace,xurl,hyperref}
\setlength{\parindent}{0pt}
\setlength{\parskip}{5pt}
\setlength{\emergencystretch}{3em}
\newcommand{\N}{\mathbb{N}}
\newcommand{\Z}{\mathbb{Z}}
\newcommand{\Q}{\mathbb{Q}}
\newcommand{\R}{\mathbb{R}}
\newcommand{\C}{\mathbb{C}}
\newcommand{\F}{\mathbb{F}}
\newcommand{\A}{\mathbb{A}}
\newcommand{\PP}{\mathbb{P}}
\newcommand{\E}{\mathbb{E}}
\newcommand{\eps}{\varepsilon}
\newcommand{\dd}{\,\mathrm{d}}
\begin{document}
\section*{UnsolvedMath ID: 6000008}
\textbf{Source problem code:} catalogue entry
\subsection{Definitions and mathematical statement}
For a nondegenerate equiaffine immersion $f:M^2\to\R^3$ with affine normal $\xi$, the affine fundamental form $h$ is the bilinear form in the decomposition
\[
D_X(f_*Y)=f_*(\nabla_XY)+h(X,Y)\xi,
\]
where $D$ is the ambient flat connection. With $D_X\xi=-f_*(SX)$, the affine mean curvature in a usual convention is $\tfrac12\operatorname{tr}S$. If ``affine minimal'' means its vanishing, the question concerns the corresponding class of immersions; the source does not give this convention explicitly.

For positive-definite $h$, local isothermal coordinates $(u,v)$ mean
\[
h=e^{2\varphi(u,v)}(\mathrm du^2+\mathrm dv^2)
\]
for a smooth real function $\varphi$. The question asks whether affine minimal surfaces admit such coordinates with respect to $h$, in analogy with Euclidean minimal surfaces.

Under this strictly local, positive-definite interpretation the supplied summary identifies an affirmative answer from the ordinary two-dimensional isothermal-coordinate theorem, which does not require minimality. For a nondegenerate indefinite form, the analogous local conformal expression uses $\mathrm du^2-\mathrm dv^2$, rather than the displayed positive-definite expression. A degenerate form is outside both interpretations.

The record does not say whether the coordinates are local or global, whether the affine form is definite, or whether additional compatibility with the affine immersion is intended. A conformal coordinate atlas is also different from one global coordinate chart. The cited local metric observation does not resolve any unstated global or stronger affine-coordinate requirement.
\subsection{Short English statement}
Do affine minimal surfaces have conformal coordinates for their affine fundamental form? The supplied answer covers ordinary local isothermal coordinates for a nondegenerate metric; global, signature, and any stronger affine requirements remain unspecified.
\subsection{Sources}
\begin{itemize}
\item Local UnsolvedMath registry, permanent ID 6000008, exact source code catalogue entry, ``Isothermal Coordinates for Affine Minimal Surfaces''; proposed by Takashi Kurose, 1998. Registry status partially solved, classification PARTIAL-PROGRESS.
\item Furuhata, Matsuzoe, and Urakawa, \emph{Open Problems in Affine Differential Geometry and Related Topics} (1998), item 3(b), scanned PDF pages 1--2, \url{https://www.jstage.jst.go.jp/article/iis/4/2/4_2_125/_pdf/-char/en}. Supplied manual OCR from page images accessed July 29, 2026. No PDF or page image was accessed here.
\item The supplied summary invokes the classical local isothermal-coordinate theorem for two-dimensional metrics and Blaschke's affine surface theory, without a precise reference. It reports the local interpretation affirmative, not a global-coordinate theorem.
\item Provenance says the historical article presents the item as open and flags current status NEEDS\_REVIEW. The supplied rights note marks the J-STAGE article CC BY 4.0; attribution is retained.
\end{itemize}
\end{document}
